\documentclass[../../main.tex]{subfiles}
\begin{document}

{\bf [解答]}

$t=10^{10}+3$として，与えられた数
\begin{align}
  A = \dfrac{10^{210}}{10^{10}+3} = \dfrac{10^{210}}{t}
\end{align}
について考える．

まず$A$の桁数について
\begin{align}
  \dfrac{10^{210}}{10^{11}}<A<\dfrac{10^{210}}{10^{10}} \\
  \therefore \,10^{199}<A<10^{200}
\end{align}
だから，$A$は$200$桁である．$\cdots$(答)

次に$1$位の数を求める．合同式の法を$10$とすると，$A$は
\begin{align}
  A
  &=\dfrac{(t-3)^{21}}{t} \\
  &=\sum_{k=1}^{21}{}_{21}C_kt^{k-1}(-3)^{21-k}-\dfrac{3^{21}}{t} \\
  &\equiv \sum_{k=1}^{21}{}_{21}C_k3^{k-1}(-3)^{21-k}-\dfrac{3^{21}}{t} \quad (\because t\equiv 3)\\
  &=\sum_{k=1}^{21}(-1)^{21-k}{}_{21}C_k3^{20}-\dfrac{3^{21}}{t} \\
  &=-3^{20}\sum_{k=1}^{21}(-1)^{k}{}_{21}C_k-\dfrac{3^{21}}{t} \label{eq:1}
\end{align}
である．

ここで二項定理より
\begin{align}
  (1+x)^{21}=\sum_{k=0}^{21}{}_{21}C_kx^k
\end{align}
だから，$x=-1$を代入して
\begin{align}
  0 &={}_{21}C_0+\sum_{k=1}^{21}{}_{21}C_k(-1)^k \\
  \therefore
  0 &=1+\sum_{k=1}^{21}{}_{21}C_k(-1)^k
\end{align}
である．これを\cref{eq:1}に代入して
\begin{align}
  A\equiv 3^{20}-\dfrac{3^{21}}{t}\tag{2}
\end{align}
である．$9 \equiv -1$だから
\begin{align}
  A = 9^{10} -\dfrac{3^{21}}{t}
  \equiv (-1)^{10} -\dfrac{3^{21}}{t}
  \equiv 1 - \dfrac{3^{21}}{t} \label{eq:2}
\end{align}
である．

$3^{21}=10460353203$を代入して
\begin{align}
  \dfrac{3^21}{t}=1.04\dots
\end{align}
だから\cref{eq:2}に代入して
\begin{align}
  A \equiv 1 - 2 \equiv 9
\end{align}
が$1$の位である．$\cdots$(答)

{\bf [別解]}

$1$の位について考える．新しく
\begin{align}
  B=\dfrac{10^{210}+3^{21}}{t}, \quad C=\dfrac{3^{21}}{t}
\end{align}
とすれば，
\begin{align}
  A = B - C
\end{align}
であるから，$B,C$の$1$の位を求めれば$A$の$1$の位が求まる．

まず$B$について，分子を因数分解すると$x=10^{10},y=3$として
\begin{align}
  10^{210}+3^{21} = x^21+y^21 = (x+y)(x^{20}+x^{19}y+\cdots xy^{19}+y^{20})
\end{align}
だから
\begin{align}
  B = x^{20}+x^{19}y+\cdots xy^{19}+y^{20} \equiv y^{20} = 3^{20} \equiv 1
\end{align}
である．ついで$B=1.04\dots$だったから，
\begin{align}
  A\equiv B-C\equiv 9
\end{align}
である．$\cdots$(答)

\newpage
\end{document}
